\appendix
\onecolumn

\section{The reconstruction procedure}
\label{app:procedure}

\paragraph{Decomposition.} The self-hosting target is decomposed into 21
components ordered by dependency. Each component declares the pinned
inputs its obligations read, the outputs a producer must supply, and the
components it depends on. The decomposition is part of the committed
material and therefore part of the identity.

\paragraph{Per-component reconstruction.} For each component in
dependency order, the harness materializes a workspace containing the
committed content only: the acceptance programs, fixtures, pinned
inputs, and the recipe. Already-reconstructed dependencies are placed in
the workspace as context, because an implementation of a higher layer
must import the lower layers it builds on. The original implementation
of the component under reconstruction is not present in the workspace at
any point.

A producer is then invoked on the workspace and its output sealed. If the
component's obligations fail, the producer is re-invoked with the failure
output, up to a per-component attempt limit. Attempt counts are recorded
and reported raw (\S\ref{app:measurements}); a reconstruction that needed
retries is not thereby disqualified, since the obligations are the
acceptance condition and not the number of attempts.

\paragraph{Whole-system judgment.} After all components are
reconstructed, the assembled tree is judged against the unchanged
whole-system identity with all 37 generated modules substituted for the
originals. This is a stronger condition than per-component acceptance:
each component passed its own obligations in isolation, and the assembled
system must additionally satisfy the system-level obligations, including
the recursive audit of the dependency chain.

\paragraph{The blindness protocol, and its limits.} Blindness means
exactly this: the harness does not place the original implementation in
the producer's workspace, and the producer is invoked with no reference
to it. It does not mean, and the protocol cannot establish, that the
producer has no prior exposure to the target. The supporting evidence is
indirect and reported in \S\ref{sec:selfhost}: measured artifact
dissimilarity, two unrelated producer families, and a held-out battery
maintained outside the committed identity. The specification states the
limitation directly: producer independence is recorded, never enforced,
because content cannot prove blindness.

\paragraph{The recursive step.} A qualifying reconstruction must also
perform the operations the system exists to perform. The step exports a
claim through the reconstruction's own command-line interface, imports it
into a fresh workspace, rebuilds its generated outputs through a
producer, and checks that the resulting identity equals the original and
that the parent verifies the child. In the generational trials, every
such operation is executed by a prior reconstruction rather than by the
shipped implementation.

The step operates on one component claim, by default the core layer, and
not on the whole system. It establishes that a reconstruction can drive
the export, rebuild and verification path. It does not establish that a
further complete system was grown.

\section{Verdicts sealed as claims}
\label{app:trialclaims}

A verdict recorded in a log is unverifiable. Each qualifying
reconstruction is therefore sealed as a claim whose own acceptance
procedure re-executes the conditions that qualified it: identity under
full substitution, the held-out battery, and the operational steps. A
claim cannot seal unless its obligations pass, so a non-qualifying trial
produces no sealed manifest. We verified this property directly during
the evidence audit for this manuscript: of seven trial directories, the
five qualifying ones carry a sealed manifest and the two non-qualifying
ones carry none.

The construction made one catch of its own before any trial sealed: it
refused a seal over a boundary whose acceptance obligations named a path
outside the claim, which was then corrected.

The regress this invites terminates deliberately. Judges are auditable by
the machinery they judge with, and the one operation that is never
automated is authorization. Adopting an obligation, and so moving an
identity, requires a human signature; every obligation discussed in this
paper carries one, recorded alongside the failure that selected it.

\section{Detailed measurements}
\label{app:measurements}

\subsection{Producer cost and timing}

Across 26 recorded whole-system reconstructions, each of 21 components:

\begin{center}\small
\begin{tabular}{@{}lrrr@{}}
\toprule
family & median wall-clock & range & per-component retries \\
\midrule
codex (\texttt{gpt-6-sol}) & 40 min & 27--211 min & 0--1 \\
claude (\texttt{sonnet})   & 186 min & 147--438 min & 0--4 \\
\bottomrule
\end{tabular}
\end{center}

The declared cost envelope for a whole-system claim is \texttt{usd =
40.0}, pinned from the project's own ledger history with headroom. The
codex family ran on a flat-rate plan and recorded zero metered spend
against that ceiling; the claude family likewise ran on a subscription
rather than metered API billing. Consequently the dollar figures in the
repository understate the market cost of these reconstructions, and we
make no cost claim from them (\S\ref{sec:discussion}).

\subsection{Implementation agreement per component}

Mean pairwise raw text similarity across the seven independent
reconstructions (the shipped implementation excluded), by component,
highest to lowest. The pattern of interest is that
agreement tracks obligation density, and that every failure witness
recorded during the trials occurred at a low-agreement component.

\begin{center}\small
\begin{tabular}{@{}lll@{}}
\toprule
component & agreement & note \\
\midrule
\texttt{\_kernel/core.py}      & 0.377 & most heavily pinned component \\
\texttt{feedback.py}           & 0.244 & \\
\texttt{\_kernel/seal.py}      & 0.238 & \\
\texttt{\_kernel/attest.py}    & 0.216 & \\
\texttt{\_kernel/recipe.py}    & 0.209 & \\
\texttt{\_kernel/identity.py}  & 0.197 & \\
\texttt{\_util.py}             & 0.182 & \\
\texttt{hooks.py}              & 0.144 & \\
\texttt{transfer.py}           & 0.140 & \\
\texttt{reference.py}          & 0.134 & \\
\texttt{\_cli/output.py}       & 0.131 & \\
\texttt{pack.py}               & 0.117 & witness: record round-trip \\
\texttt{\_kernel/run.py}       & 0.116 & witness: sandbox floor \\
\texttt{attest.py}             & 0.114 & \\
\texttt{\_cli/handlers.py}     & 0.111 & \\
\texttt{assess.py}             & 0.110 & \\
\texttt{render.py}             & 0.104 & \\
\texttt{record.py}             & 0.102 & \\
\texttt{authoring.py}          & 0.099 & \\
\texttt{\_kernel/crosscheck.py}& 0.090 & \\
\texttt{\_cli/parser.py}       & 0.088 & \\
\texttt{\_cli/views.py}        & 0.084 & \\
\texttt{reuse.py}              & 0.075 & \\
\texttt{launcher.py}           & 0.067 & \\
\texttt{\_cli/statusview.py}   & 0.066 & \\
\texttt{\_kernel/build.py}     & 0.064 & \\
\texttt{kernel.py}             & 0.061 & \\
\texttt{registry.py}           & 0.060 & witness: staging, resolution \\
\texttt{\_cli/report.py}       & 0.057 & \\
\texttt{heldout.py}            & 0.049 & \\
\texttt{\_cli/verbs.py}        & 0.046 & \\
\texttt{\_cli/dispatch.py}     & 0.024 & least constrained component \\
\bottomrule
\end{tabular}
\end{center}

Package initialization files and the top-level entry point are omitted;
they agree at 0.58--0.76 and contain almost nothing to write. We state
the relationship between agreement and obligation density as an
observation rather than a tested correlation: the components are not
independent observations, and obligation density was not assigned
experimentally.

\subsection{Diversity instrument}

The similarity figures in \S\ref{sec:selfhost} were produced by an
instrument fixed before the measurements reported here were taken, and
validated by reproducing an earlier reading of the project's own corpus.
On the raw channel it recovers that reading to within $0.003$ on seven of
eight per-reconstruction scores, with the eighth high by $0.0070$, and
reproduces the count of 32 of 37 components structurally distinct in
every reconstruction. It reports two channels: raw text, and text with
comments, docstrings and layout normalized away via syntax-tree
unparsing. The normalized channel runs 0.05 to 0.09 higher, because a
substantial part of what separates two reconstructions as text is prose.

\section{External arm detail}
\label{app:jq}

\subsection{Materials}

The corpus is the official \texttt{jq} 1.8.1 test suite, 522 parseable
cases, frozen with checksums. All three implementations return the same
verdict on 373 of them --- 332 passed by all three, 41 failed
identically by all three --- and differ on the remaining 149. The oracles are \texttt{jq} 1.8.1,
\texttt{gojq} 0.12.19 and \texttt{jaq} 3.1.1, also frozen. The oracles
adjudicate disagreement and are never used as acceptance criteria, so
that no obligation is defined by running a reference implementation
inside the gate.

The initial specification $S_0$ is the evaluation-core slice: the 465
cases exercising expression semantics, excluding regular expressions,
dates, input/output, modules and command-line flags, and excluding
approximately 43 cases that fail even for the reference implementation
under a plain runner. The carve rule is "the case drives only the
evaluator", applied by a tagger script and reviewed once.

\subsection{Obligation inventories}

$S_1$ adds 16 computed and metamorphic laws evaluated over inputs seeded
from entropy at evaluation time: indexing and length relations computed
independently by the checker; composition against sequencing; idempotence
of sorting; round-tripping of serialization; and concatenation length
additivity. Each was validated on three fresh seeds against the
reference before adoption.

$S_2$ adds 28 error-condition laws requiring that generated ill-typed
inputs raise, covering wrong-typed arguments to entry conversion,
transposition, joining and character conversion; malformed path arrays
given to path deletion and assignment; non-path expressions given to path
extraction; mismatched operand types for containment and indexing; and
arithmetic on incompatible types. Each condition was first verified
unanimous across all three implementations, and the battery validated on
five fresh seeds. A sanity guard accompanies the battery so that an
implementation refusing all inputs cannot satisfy it.

Three drafted laws were rejected before adoption because measurement
contradicted them. The reference's refusals for list reversal,
transposition and vacuous quantification proved value-dependent rather
than type-dependent; reversing the number $0$ yields the empty list while
reversing $5$ raises, because the definition routes through a length that
is zero in the first case. One candidate was abandoned entirely, its
refusal condition not being expressible as a property of types.

\subsection{Probe battery and per-sample results}

The differential battery is 1{,}395 deterministic seeded probes over
program and input pairs that the corpus does not contain, executed
against four implementations. The human-divergence surface is the 340
distinct probes on which \texttt{gojq} or \texttt{jaq} disagrees with
the reference, giving a base rate of $340/1395 = 24.4\%$; it was frozen independently of
any reconstruction.

\begin{center}\small
\begin{tabular}{@{}llrr@{}}
\toprule
boundary & sample & divergent probes & on human surface \\
\midrule
$S_1$ & codex 1 & 79 & 47 (59\%) \\
$S_1$ & codex 2 & 52 & 43 (83\%) \\
$S_1$ & codex 3 & 66 & 44 (67\%) \\
$S_1$ & pooled (distinct) & 92 & 51 (55\%) \\
$S_1$ & claude 1 & 34 & 26 (76\%) \\
\bottomrule
\end{tabular}
\end{center}

Of the 92 pooled codex divergences, 43 were flagged by all three samples,
19 by two and 30 by one, which indicates the battery measures structure
rather than noise. Of the 43 flagged by all three, 40 (93\%) lie on the
human-divergence surface.

The claude sample is 3{,}263 lines against the codex samples' 869, a
ratio of 3.8, which independently reproduces the family size signature
measured on the self-hosting corpus. It is also qualitatively different:
17 of its 34 divergences are cases where both implementations succeed
with different values, where the codex samples' divergences were
dominated by excess permissiveness.

\subsection{Divergence classes by family}

The class decomposition that explains the negative cross-family result in
\S\ref{sec:jq}:

\begin{center}\small
\begin{tabular}{@{}llrrr@{}}
\toprule
family & boundary & permissive & value split & sample error \\
\midrule
codex  & $S_1$ & 35.3 & 19.3 & 10.0 \\
codex  & $S_2$ & 27.5 & 13.0 &  8.0 \\
claude & $S_1$ &  6   & 17   & 11   \\
claude & $S_2$ &  5   & 20   &  9   \\
\bottomrule
\end{tabular}
\end{center}

Excess permissiveness is 54\% of the codex family's divergence and 18\%
of the claude family's. The $S_2$ repair targeted permissiveness, and
permissiveness fell in both families. The aggregate moved only for the
family in which that class dominates.

\subsection{The unresolved election}

Of the programs diverging consistently at both boundaries, 19 reduce to
four classes: six differing only in error message text, which project
precedent leaves unpinned; nine with a single reimplementation as the
outlier, admitting a mechanical consensus resolution; three constituting
a genuine election, where indexing an object by a string key returns null
in the reference and raises in both reimplementations; and one liveness
divergence recorded as a finding rather than pinned. The three-case
election is unresolved. Electing the reference requires members to
replicate a permissive edge two independent implementations treat as an
error; electing the majority produces a specification that refuses the
artifact it was recovered from.

\section{Reproduction}
\label{app:repro}

The repository contains the protocol for an independent operator to
re-earn a sealed trial, together with the signed-receipt format and a
continuous-integration workflow that checks a returned receipt against
the bundle's own public key. Re-earning a sealed trial claim requires the
checkout, the retained reconstruction lineage, and one audit command; it
executes the acceptance obligations and does not require a producer or
any model spend.

Reproducing a reconstruction from scratch requires a producer. The
expected cost is given in \S\ref{app:measurements}; because both families
were run on flat-rate plans, an operator billed per token should expect a
materially higher figure, and the declared envelope of \texttt{usd =
40.0} per whole-system claim is the project's own stated ceiling rather
than a measured market price.

No external operator has returned a receipt. The two records in the
repository named for this condition are same-host dry runs of the
instrument, retained to show the measurement path works, and each still
carries a placeholder note in place of an operator's signature.

\twocolumn
